\section{Case Study: A Voice Coworker That Uses a Computer}\label{sec:case}

The benchmarks measure turns and tool calls one at a time. Our launch demonstration asks what interaction mode
is for: a conversation with an agent that does real work on a computer while the conversation goes on, the way
one works with a colleague, asking for one thing, chatting, asking for the next, and approving what changes
something. We recorded it in one take, with no edits to what was said or done.

\paragraph{The set-up.} The session's background is Claude Opus~5.5 as an orchestrator, with three agents
registered as the operator's webhook backends (the settings \texttt{tasks.orchestrator} and
\texttt{interaction.task\_backends}):
\begin{itemize}
\item \textbf{browser}: a Claude Opus~5.5 tool loop driving a visible Chromium window through Playwright (open a
page, follow a link, read the page while scrolling through it, highlight what it found);
\item \textbf{terminal}: a Claude Opus~5.5 tool loop that types commands into a visible terminal (a tmux session in
an xterm) and reads their output;
\item \textbf{writer}: a Claude Opus~5.5 tool loop that types a draft into LibreOffice Writer and, after the caller
approves, saves it.
\end{itemize}
The agents run on a sandboxed Linux desktop (a container with a 1920$\times$1080 virtual screen, no access to the
host's files or local services, and internet access), whose screen is recorded. Their actions are paced for a
person to follow: keys are typed one at a time (25--60\,ms apart in the terminal, 16\,ms in Writer), and the
browser's pointer glides to a link before it clicks; the work itself is real. Each agent reports progress,
questions and results with the signed webhook protocol of Section~\ref{sec:system}; the save is a side effect, so
it is read back to the caller and needs a spoken yes. The caller's voice is synthesised (ElevenLabs) before the
call and each line is said when the call reaches the moment it is written for, the ``yes'' only after the
read-back has been heard and the agent has stopped talking. Speech recognition and the agent's voice are
Deepgram (Aura-2 ``juno''); the front model is Qwen3.5-2B on the same GPU. The take ran on a build of the gateway
with the fixes described below (PR~\#49 on the project's repository), not on release 1.1.0, and with
\texttt{failure.stall\_after\_ms} raised from 2.5 to 8\,s, an operator setting for a cloud planner that takes a few
seconds before its first word.

\paragraph{What happens.} Times are seconds from the start of the call; reply latencies are measured on the
recording from the caller's last voiced sound.
\begin{itemize}
\item \textbf{1.1--6.2\,s.} ``Morning, Ava! Can you pull up the GitHub release for WaaV one point one and check it all
looks right?'' The agent acknowledges (``Let me look into that.'') 1.8\,s after the caller stops; the orchestrator
starts the browser at 10.3\,s and says what it is doing (``I'm pulling up the version one point one release on
GitHub now to check it over.'').
\item \textbf{10.3--26.3\,s.} The browser opens the release page, scrolls through it and highlights the ``Latest''
label and the release note on interaction mode. Its result (published, marked Latest, interaction mode in the
notes, all three files) is told at 27.1\,s.
\item \textbf{37.0--45.3\,s.} The caller makes conversation: ``Thanks. Remind me, how many speech vendors do we
support now?'' The answer comes from the orchestrator, from the facts it was given: ``Forty-four speech vendors.
That's thirty-one for speech-to-text, thirty-four for text-to-speech and eleven for speech-to-speech \ldots''
\item \textbf{54.1--77.0\,s.} ``Great. Can you also get the download numbers in the terminal?'' The terminal starts at
60.3\,s, types a \texttt{curl \textbar\ jq} query against the GitHub API, and finishes at 77.0\,s: the x86-64 archive
had 0 downloads, the arm64 archive 12 and the checksum file 14.
\item \textbf{92.0--124.2\,s.} ``Nice. Can you draft the launch announcement now, with those numbers in it?'' The
orchestrator passes the numbers found earlier to the writer (started at 102.1\,s), which types the announcement
into Writer on screen, title, sections, bullet lists and an ``Early numbers'' section, while the orchestrator says
what it is doing.
\item \textbf{124.6--148.5\,s.} The save is read back: ``To confirm: save the document, Bud WaaV 1.1.0 Is Out. Do you
want me to do that?''. The caller's ``Yes, go ahead.'' is accepted (``Okay.''), and the document is saved at
148.5\,s and told at 149.7\,s.
\item \textbf{158.0--161.6\,s.} ``Perfect, that's everything. Thanks, Ava!'' --- ``You're welcome.'' The call lasts
168\,s.
\end{itemize}
Every request was acknowledged 1.5--1.8\,s after the caller stopped speaking; each piece of work started on the
right agent, ran while the conversation went on, and was told when it finished.

\paragraph{What went wrong in this take.} Two defects are visible and are reported as they happened. The front
model's spoken update of the terminal's result was garbled (``one point one point zero for the x86 build, twelve for
ARM''); the orchestrator corrected it on its own 25\,s later (``Quick fix on what you heard before: the x86 build
actually has zero downloads, ARM has twelve, and the checksum file has fourteen. That's twenty-six in total.''). And
the caller's yes came 9.6\,s after the question, because our caller script waits for the line to go quiet; the
gateway, having heard nothing for 8\,s, asked again in short form (``Do you want me to save the document?'') just
as the yes began, and the yes was then accepted.

\paragraph{What the takes before it found.} We recorded 27 takes while the host was shared with unrelated GPU jobs
and on Wi-Fi. The earlier ones exposed defects in the gateway that the takes were then used to fix (PR~\#49): task
results lost because the narrator's length limits disagreed or a version number ended its sentence; an update
spoken under the wrong task; a sentence cut by a question replayed whole; holding lines that repeated the line
before; after a task had finished, new requests read as follow-ups about it (a 2B front chose \texttt{follow\_up}
for 10 of 10 such turns offline; with a check that names the finished task's agent, 130 of 130 of a gate set were
settled right); a rejected synthesis flush from the speech vendor that left every later line unfinished and played
an apology every 13\,s for six minutes; a request split at ``the'' and read as a duplicate; a late transcript that
turned a clear yes into a weak one; and an unanswered offer that held the results back for minutes. Other takes
failed on the network (the speech vendor's connection timed out at set-up) or on the shared GPU (the front model's
decisions missed their budget). We chose the demonstrated take for being complete, not for being typical, and the
remaining failures are in the project's issue list.
